\section{Introduction}
\label{iv:introduccion}
\begingroup
\fontsize{10.5}{13.2}\selectfont
\setlength{\parskip}{1pt plus .5pt minus .25pt}

How can a single discrete structure determine coordinates of the
continuum, support an exceptional realisation and preserve a duality
on passing to dimensional representations? The question requires more
than collecting values, dimensions and groups: it requires the construction
of intermediate objects and the maps that transmit their information.
This article develops that connection in Triadic Modular Holography (HMT).
Its organising result is a twofold extension of the genealogical state:
incidence leads to a vertex operator algebra isomorphic to
\(V^\natural\), while phase, action and channel representations
determine a compact duality and a screen morphism. Coordination of
these extensions then specifies the algebraic domain and transport
conditions for an M-theory realisation.

The construction begins before real-valued evaluation.
\emph{All Possible Paths} (APP) denotes the paths on the support
\(X_9=(\mathbb Z/9\mathbb Z)^2\). Its Cayley graph uses the additive
law and the cardinal-direction generators \(\pm(1,0),\pm(0,1)\).
Sum and product are evaluated on the same positions; residual reduction
retains the quotient of each evaluation. The additive and multiplicative
sheets are therefore two evaluations on a common support.
The signature \(\mathrm{TRIT}=\{-1,0,+1\}\) distinguishes transport
orientation and the hyperbolic, parabolic and elliptic regimes of its
quadratic realisation. The \emph{Topological Phase Kernel} (TPK) is the
effective composition of evaluation selection, cardinal-direction
transport and register update that extends compatible states. These
states retain positions, sheets, orientation, residues, quotients,
carries, prefixes, boundary data and memory. Selection does not replace
displacement, and phase return does not erase the history traversed.

% BEGIN R27_FRAMING_INTRO_EN
The operations connecting arithmetic support with incidence admit
explicit intermediate proofs. Changing the modulus, repeating the support
at fixed modulus, and varying the dimension have their own sum--product
defect laws. Compression into three classes produces an index-two
sublattice and a quadratic unit of norm one. Within the emitter, the
decimal alphabet, half-emission junctions, and the contrast between
histogram and ordered memory preserve the distinctions between value,
incidence, and trajectory. The multiplicative observable is examined
before regional selection.

% BEGIN R28_FRAMING_INTRO_EN
The emitter's labelled multigraph and its phase quotient place these
realizations on explicit incidences. The five minimal eight-edge closures
cover the APP, propagation, autoscale, and individual closure-region
anchors. Minimum boundary occupancy and the parity of memory readers
further specify representative selection and the transformation of
records under route reversal.
% END R28_FRAMING_INTRO_EN

The first result is the joint organisation of the discrete structure
of the continuum. In Section~\ref{iv:continuo-conjunto}, a single tree
of executions determines cylinder measure and sheet transport,
oriented balance, ternary incidence, the memory complex and its terminal
assembly with determinant line. Maps are defined on generators, and
their compatibility with history restrictions and modular reductions
is proved. The measure uses compatible positive weights on a finitely
branching tree without terminal leaves; the Gram identity retains the
terminal survival remainder; the determinant construction retains its
degrees and bases, together with the additional hypotheses required by
a scalar evaluation. The five constructions thus retain a common origin
and the particular conditions of each operation. They are not five
domains selected after the coordinates have been produced.

Sections of compatible cylinders determine the numerical coordinates of
the correlated generation of \(\pi,\varphi,e,\alpha\) under the
nesting, contraction and boundary conditions of their Archimedean
realisation. The exceptional projection instead retains the incidence
data realised as codes, lattices and graded algebras. The areal and
volumetric sheets and their measures, developed in
Section~\ref{sec:medida-retorno-hojas}, and the vacancies associated
with refinement, in Section~\ref{vac-capacidad-trit}, enter the
construction of \(\alpha\) and the exceptional extension together
with the regional registers. Electron mass uses a different composition
of these common operations; its evaluation is not the subject of this
article.

The chart grouping two trits into a nonadic digit has an inverse and
a transported addition with carry; the projection of the ambient cube
of \(729\) words onto its first two coordinates has uniform nine-element
fibres. The quadratic-correlation proof then
determines the identity of the six-position chart. At the boundary,
saturation leaves eight possibilities, dual neutrality leaves two,
and orientation selects the fixed representative. Closure controls retain
the endpoints of the telescoping identity. 

The exceptional connection has a concrete antecedent in the
dodecaphasic register \(K\). Its twelve components arise from the
signed register and the TPK operations. Rational encoding preserves
the blocks when base, length, phase origin and orientation are fixed;
Hadamard inversion recovers the corresponding signed coordinate.
The supports of the register enter a marked flag of the Witt design.
The ternary code constructs the gluing of the Niemeier lattice with
root system \(A_2^{12}\); the flag and origin, together with the
specified metric and positive chamber, select the rootless neighbour
identified with the Leech lattice. Each operation transmits a datum
needed by the next object.

% BEGIN S27_FRAMING
Section~\ref{star27:comparacion} specifies the passage from regional
registers to incidence through a comparison of three stars.
The intersections of \(P_\pi\) with \(H_e\), \(H_\varphi\) and
\(H_{\rm APP}\) determine the four-subsets \(B_K\), \(B_\varphi\)
and \(B_{\rm APP}\); the sign partition of the pre-carry cochain
determines the polarisation of their completions. The census of
all \(495\) four-subsets locates these configurations within the
complete design, and the stabiliser of the ordered four-support
frame characterises its rigidity within \(G_W\). This rigidity
concerns the action on positions; the numerical determination of
\(K\) retains its integer construction, phase origin and previously
fixed orientation.

Section~\ref{carry27:seccion} develops the integer component accompanying
the residue reading. Euclidean division produces the residues \(q,a\)
and quotients \(c,h\), whose composition retains the memory of multiples
of three. The cocycle, polynomial representations over \(\mathbb F_3\)
and transport identities specify which data pass through each change
of register. Their connection to compatible cylinders retains the
boundary bounds and the fibres of the regional reading. The star
comparison and carry law thus have specified incidence and numerical
roles in the connection that this article extends to lattices and
algebras.

% END S27_FRAMING


The lattice algebra is constructed on this lattice with oscillators,
central extension and vertex product. The order-two extension of
Frenkel--Lepowsky--Meurman and the order-three extension associated
with the ternary isometry realise \(V^\natural\), once the hypotheses
of the recognition theorems used have been verified
\cite{flm,chenlamshimakura2016,abelamyamada2017}.
The chosen isomorphisms between the presentations preserve the vacuum,
conformal vector, grading and products, and transport automorphisms
and traces. The Monster is the automorphism group of this algebra;
its action has domain \(V^\natural\), not the digits of \(\alpha\).
Borcherds's Moonshine theorem determines the corresponding graded
traces~\cite{borcherds}. Lattice construction, algebraic extension and
recognition of these traces are distinct stages of one chain of maps.

Upon reaching the level-three
traces, formal factorization of the rational identity determines
divisibility at every degree, with its optimal exponent.
% END R27_FRAMING_INTRO_EN

The register \(K\) also has a directional realisation.
In the specified representation of \(A_5\) on the twelve positions,
its projection onto the three-dimensional sector determines a vector
\(u_K\ne0\) in \(V_{11}=\mathbf1^\perp\).
The nonzero vector determines the line \(\ell_K\), and its exact squared
norm normalizes the rank-one term of the projector onto
\(V_{10}=V_{11}\cap u_K^\perp\), with specified
kernel and covariance. The reduction \(12\to11\) removes the uniform
mode; the reduction \(11\to10\) removes the direction selected by
the complete register. Its cyclic orbit is also an invertible reference
for identifying operators through their responses.
Section~\ref{sec:acoplamiento-recuperacion} establishes the stability
bounds and bilateral balance: two channels constructed from isometries allow the state
to be recovered; the unitary extension includes incoming memory,
and preservation of an observable requires the commutation condition
in its statement.

The Hilbert-space realisation supplies the operator-theoretic support
for duality. At each depth, phases label an orthonormal basis and the
shift and character operators generate its full matrix algebra.
Their Weyl relation represents the central extension of translations
with the sign of the oriented area form fixed. Inclusion through the
identity of the entire new fibre preserves the unit and normalised
trace: the limit is the nonadic uniformly hyperfinite algebra, and
its tracial closure is a factor of type \(II_1\). Fixing one
continuation instead gives a corner inclusion and a different closure.
The refinement map therefore determines the type of realisation,
not merely its dimension.

The action section and angular coordinates fix the radius that
specialises the compact duality. The determinantal norm, the coefficients
of the action character and normalised regional renewal determine the
positive sections before and after return. The angular charts retain
orientation and winding number. Their composition defines a labelled
ellipse whose semiaxis product preserves action and whose ratio
determines a unique dimensionless radius. Positivity conditions place
the generated coordinates in the elliptic chart; semiaxis exchange
produces the reciprocal radius while preserving the common action scale.

On \(\mathbb Z^2\times\mathbb R_{>0}\), integer-coordinate exchange
and radius inversion form an involution preserving the compact
quadratic form. The diagonal Hamiltonian is self-adjoint, nonnegative
and has compact resolvent; coordinate permutation unitarily conjugates
it to the reciprocal-radius Hamiltonian, including the operator and
quadratic-form domains. A change of residual section also transports
the relation \(L_9\) that preserves the integer. The minimum energy
of each class is well defined and is evaluated using two integer
candidates. Normalisation by the string scale expresses the same
duality in dimensional radii. On the imaginary axis of the upper
half-plane, modular inversion is related to it by an explicit
semiconjugacy of parameters. This map does not identify the exchange
operator with the automorphisms of orbifold sectors.

The dimensional screens coordinate the \(K\)-direction, incidence
and compact sector. Puncturing the ternary code produces a length-eleven
code with distance five and a perfect partition. The orthogonal
reduction \(12\to11\to10\) and the isotropic lattice reduction
\(26\to24\) retain their own presentations, relations and kernels.
The morphism \(\mathcal R_M^{\mathrm{alg}}\) identifies their quotient
presentations and intertwines duality; the graph of the projection
retains the coordinate in \(V_{11}\) together with its image in
\(V_{10}\). The same operations restrict to the collection of the two
radii determined by the ellipse.

The dynamical extension begins with trajectory composition, not with
the choice of fields. Retaining the preceding datum at the boundary
makes changes of direction and arithmetic type an additive action
cocycle. The worldsheet and Polyakov action then specify a variational
realisation. The map to the M-theory realisation receives the screens
in an eleven-dimensional codomain equipped with a metric, three-form,
gravitino, grading, action and constraints. Under the intertwining and
domain-transport hypotheses, the supercharge theorem preserves
\(Q^2=H\) on the image of the isometry. Complete physical identification
additionally requires tensions, oscillators and amplitudes, anomaly
control and the limits stated in
Section~\ref{iv:condiciones-realizacion-fisica}.

This hierarchy determines the development: kernel and joint continuum;
coinductive generation, register \(K\) and the two constructions of
\(\alpha\); incidence and exceptional algebra; Hilbert-space realisation,
action and radius; duality, screens and dynamical transport.
Section~\ref{sec:k-registro} determines the register used in
\ref{exc:k-direccion}; the realisation in \ref{iv:hilbert} supports
the operators in \ref{iv:dualidad}; and \ref{iv:pantallas} composes
the reductions in the morphism received by the eleven-dimensional
realisation. The argument thus connects structure, value, representation
and equation through their effective maps.

\par\endgroup
